\documentclass[11pt]{article}

\usepackage[T1]{fontenc}
\usepackage[utf8]{inputenc}
\usepackage{lmodern}
\usepackage[margin=1in]{geometry}
\usepackage{microtype}
\usepackage{amsmath,amssymb,amsthm,mathtools}
\usepackage{booktabs,longtable,tabularx,array}
\usepackage{enumitem}
\usepackage{xcolor}
\usepackage{hyperref}

\hypersetup{
  colorlinks=true,
  linkcolor=blue!55!black,
  citecolor=green!45!black,
  urlcolor=blue!65!black,
  pdftitle={Finite-Regulator Source Insertions and Obstruction Tests},
  pdfauthor={YesterdaysLemon; AI-assisted open research ledger}
}

\setlist{nosep}
\setlength{\parindent}{0pt}
\setlength{\parskip}{0.55em}
\setlength{\emergencystretch}{3em}

\newtheorem{theorem}{Theorem}[section]
\newtheorem{proposition}[theorem]{Proposition}
\newtheorem{lemma}[theorem]{Lemma}
\newtheorem{corollary}[theorem]{Corollary}
\theoremstyle{definition}
\newtheorem{definition}[theorem]{Definition}
\newtheorem{remark}[theorem]{Remark}

\newcommand{\YM}{Yang--Mills}
\newcommand{\RG}{\mathrm{RG}}
\newcommand{\E}{\mathbb E}
\newcommand{\dd}{\mathrm d}
\newcommand{\Hrc}{(\mathrm H_{\mathrm{rc}})}
\newcommand{\HJ}{(\mathrm H_J^{\mathrm{conn}})}
\newcommand{\claim}[1]{\texttt{#1}}
\newcommand{\repo}{\href{https://github.com/YesterdaysLemon/yang-mills-research}
  {yang-mills-research}}

\title{\textbf{Finite-Regulator Source Insertions and Obstruction Tests}\\
\large in a Gauge-Covariant \(SU(2)\) Renormalization Step\\[0.5em]
\normalsize An exploratory research handoff}
\author{YesterdaysLemon\\
\small assembled from an AI-assisted open research ledger}
\date{23 July 2026}

\begin{document}
\maketitle

\begin{center}
\fcolorbox{red!60!black}{red!4}{
\begin{minipage}{0.91\textwidth}
\textbf{Status.}
This paper is not a solution of the \YM{} existence and mass gap problem.
It proves no continuum limit, no infinite-volume quantum field theory, and
no mass gap.  It reports finite-regulator algebra, conditional one-step
constructions, and obstruction tests from an exploratory public repository.
The repository-level evidence status is E0; the strongest individual
auxiliary claims are internally checked at E2 and have no independent human
review.
\end{minipage}}
\end{center}

\begin{abstract}
We record the outcome of an exploratory attempt to insert a local
plaquette-energy source into a gauge-covariant, Ba{\l}aban-type
renormalization step for four-dimensional \(SU(2)\) Wilson lattice gauge
theory.  On one selected finite-regulator small-field branch, the work
derives exact source-jet identities, transports one marked plaquette through
the finite \(\RG\)-II hard-core and cluster algebra, and isolates explicit
conditional hypotheses under which completed external-field derivatives
would be uniformly local.  Several natural repairs of the remaining
physical-\(U\) pullback fail: independent-link complex collars shrink like
the square of the lattice scale; diagonal scale balancing restores the same
loss through Cauchy estimates; and separate absolute values destroy possible
within-column \(U/J\) Ward cancellation.

The final checkpoint has two parts.  First, differentiation of the
variational background equations gives an exact affine range description
and elliptic factorization for the source kernel.  Second, a
source-normalized finite periodic Cartan model has an exact cochain
identity.  A localized coarse bond has nonzero coarse curl, and discrete
Stokes forces raw \(s^{-1}\) and gradient/curl \(s^{-2}\) lower bounds for
every right inverse in that model.  The retained primary sources do
\emph{not} establish that this finite model embeds into their full
multiscale \(H_1\) domain or Hessian-selected range.  That embedding,
completed weak-dual control, nonzero-source polymer analyticity, large
fields, iteration, regulator removal, reconstruction, and the infrared mass
gap all remain open.  The purpose of this paper is to preserve the exact
discoveries and failure modes in a form suitable for future researchers and
future language-model agents.
\end{abstract}

\tableofcontents

\section{Status, target, and interpretation}

\subsection{The actual problem}

The Clay problem asks, for every compact simple gauge group \(G\), for a
nontrivial pure \YM{} quantum field theory on four-dimensional space,
axiomatic control at least at the level described by Jaffe and Witten, the
required short-distance \YM{} behavior, and a reconstructed Hamiltonian
whose spectrum has a finite positive gap above the vacuum
\cite{jaffewitten}.  A finite lattice, a formal path integral, perturbation
theory, numerical glueball data, or a fixed-coupling strong-coupling gap is
not this theorem.

The working scaffold in the repository is
\[
\begin{gathered}
\text{Wilson regulator}
\longrightarrow
\text{gauge-covariant \(\RG\)}
\longrightarrow
\text{local observables}
\\[-0.2em]
\longrightarrow
\text{OS reconstruction}
\longrightarrow
\text{infrared bridge}
\longrightarrow
\text{spectral gap}.
\end{gathered}
\tag{1.1}
\]
Wilson's compact regulator supplies exact finite-dimensional gauge
invariance \cite{wilson1974}.  Osterwalder--Schrader reconstruction,
Osterwalder--Seiler lattice positivity, and L\"uscher's transfer matrix
describe the desired positivity and spectral scaffolding
\cite{os1973,os1975,osterwalderseiler,luscher}.  Ba{\l}aban's papers provide
the principal ultraviolet \(\RG\) input used here
\cite{balaban1984propagatorsII,balaban1985averaging,
balaban1985gaugefixing,balaban1985backgroundpropagators,
balaban1985variational,balaban1987rgi,balaban1988rgii}.

These inputs do not supply the arrows in (1.1).  In particular, if \(a\) is
the lattice spacing and \(T_a\) has leading spectral values
\(\lambda_0,\lambda_1\), then
\[
m_{\mathrm{lat}}(a)=-\log\frac{\lambda_1}{\lambda_0},
\qquad
\Delta_{\mathrm{phys}}(a)=\frac{m_{\mathrm{lat}}(a)}{a}.
\tag{1.2}
\]
A finite nonzero continuum gap would require
\[
m_{\mathrm{lat}}(a)=a\Delta+o(a).
\tag{1.3}
\]
A fixed positive gap in lattice units normally corresponds to a correlation
length of order one lattice spacing and does not establish (1.3).

\subsection{Evidence vocabulary}

The public ledger \repo{} separates project status from the status of an
individual auxiliary statement.  Its current interpretation is:
\begin{itemize}
  \item \textbf{E0 (project):} exploratory; no candidate solution wording is
  permitted.
  \item \textbf{E1 (individual):} a derivation or source transcription with
  incomplete independent checks.
  \item \textbf{E2 (individual):} exact finite algebra or a primary-source
  consequence with executable regressions and explicit objections.
\end{itemize}
No E2 label certifies deep mathematical truth.  The tests check identities,
finite countermodels, links, and governance boundaries.  No independent
human expert has reviewed the chain.

\subsection{Scope of this report}

The developed material concerns:
\begin{itemize}
  \item finite four-dimensional periodic regulators;
  \item primarily the group \(SU(2)\);
  \item one gauge-covariant block step;
  \item a selected near-identity small-field branch;
  \item one localized plaquette mark and predominantly the first source
  derivative at \(t=0\);
  \item conditional external-\((U,J)\) analytic tubes.
\end{itemize}
It does not identify the selected coordinate law with the unrestricted raw
block transform.  It does not construct a nonzero-source polymer disk, treat
the complete large-field operation, iterate the \(\RG\), remove either
regulator, restore all Euclidean symmetries, reconstruct a continuum
Hilbert space, or prove infrared decay.

\section{Finite source identities and the selected branch}

\subsection{Finite positive kernels}

Let \((\Omega,\nu)\) be a finite measure space, let \(w>0\) be integrable,
and let \(O\) be a bounded real observable.  Write
\[
Z(t)=\int_\Omega e^{-tO(\omega)}w(\omega)\,\dd\nu(\omega),
\qquad
\mathcal Z(t)=\frac{Z(t)}{Z(0)}.
\tag{2.1}
\]
The following basic statement is independent of gauge theory.

\begin{proposition}[Exact finite source jets]
\label{prop:source-jets}
For (2.1),
\[
\left.\partial_t\log Z(t)\right|_{t=0}=-\E(O),
\qquad
\left.\partial_t^2\log Z(t)\right|_{t=0}
=\operatorname{Var}(O).
\tag{2.2}
\]
Mixed derivatives give covariances.  The normalized ratio \(\mathcal Z\) is
unchanged by any positive, source-independent rescaling of \(w\).
Furthermore, a centered bounded observable gives an explicit zero-free
scalar disk; one convenient bound used in the ledger is
\[
|z|<\frac{2\log 2}{\operatorname{diam}O}.
\tag{2.3}
\]
\end{proposition}

This proposition gives exact finite-volume cumulants and projective
invariance.  It does not give a polymer expansion at \(t\ne0\), uniformity
in a regulator, or a continuum measure.

\subsection{Wilson fields and a local mark}

On a finite torus \(\Lambda\), an \(SU(2)\) bond field assigns
\(U(b)\in SU(2)\) to each oriented nearest-neighbor bond, with
\(U(-b)=U(b)^{-1}\).  For an oriented plaquette \(p\), let \(U(p)\) be the
ordered product around its boundary.  A representative local
plaquette-energy insertion is
\[
O_p(U)=1-\frac12\operatorname{Re}\operatorname{Tr}U(p).
\tag{2.4}
\]
The Wilson action is a sum of such terms.  A profile source would replace
the action by \(S(U)+\sum_p t_pO_p(U)\).

The research program asks for a one-step formula in which derivatives in
\(t_p\) are represented by gauge-invariant, exponentially localized
coarse operators and polymer remainders with volume-independent constants.
The first source derivative is represented algebraically by a dual number
\(\epsilon\), \(\epsilon^2=0\).  This keeps exactly one marked occurrence
through products without pretending that a common nonzero source disk has
already been proved.

\subsection{Selected coordinate law}

The Ba{\l}aban variational construction provides a local analytic
minimizing background on an admitted small-field patch
\cite{balaban1985variational}.  The repository combines that construction
with the \(\RG\)-I chart covariance
\cite{balaban1987rgi,balaban1985variational}, fixes one near-identity
branch, retains the chosen gauge chart and cutoff, and normalizes the
resulting finite integrand to a probability law.  Within one relatively
compact patch, analytic tests and scalar source ratios are holomorphic, and
real coarse gauge transformations act by push-forward.

\begin{remark}[The selected/raw distinction]
The selected branch is an operational coordinate law.  No theorem in the
ledger identifies it with the unrestricted group-delta block transform.
All source results below inherit this restriction unless explicitly stated
otherwise.
\end{remark}

\section{One marked plaquette through the finite \(\RG\) algebra}

\subsection{Distinguished-vertex Ursell algebra}

Consider a finite hard-core polymer gas with activities \(z(X)\), and add a
single marked activity \(\dot z_p(X)\epsilon\).  Differentiating the
connected logarithm selects exactly one distinguished occurrence.  The
ledger checks the connected-graph/Ursell coefficient including repeated
labels and the \(1/n!\) convention.  Non-root components have zero marked
derivative, so the derivative of a product selects the unique component
containing the mark.

This is elementary but important: expanding a derivative over all later
placements manufactures terms that are absent from the exact dual-number
product.  The mark must be routed through the actual factorization before
absolute values are taken.

\subsection{Conditioning and the whole-integrand contour}

The marked weakening integrand contains conditioning factors, Gaussian
relative-field factors, cutoffs, and later scale reorganizations.  The
repository uses one Cauchy contour for the \emph{whole} integrand.  Separate
majorants for separately differentiated factors do not bound the product
unless a common holomorphic domain and a joint integrable majorant are
available.

The source-faithful routing calculation places the local mark on the
conditional interior field \(B\), not on a shifted variable chosen only for
convenience.  For the relative Gaussian factor, the mark is independent of
the relative variable, so its additional Gaussian moment cost is exactly
zero.  This closes a finite-term routing issue; it is not a proof of a
nonzero-source gas.

\subsection{Marked resummation and a fixed-gas first jet}

After the Section-2 reorganization in Ba{\l}aban's \(\RG\)-II cluster
construction \cite{balaban1988rgii}, the ledger obtains one decorated marked
activity and a unique marked-component factorization.  Under its explicit
smallness hierarchy, doubled marked-amplitude refinement, and a separate
ordinary Koteck\'y--Preiss ceiling \cite{koteckypreiss}, the marked seed
can be resummed with positive bounds.

In the notation of the ledger, a representative bound is
\[
|W_p^{\mathrm{post}}(Z)|
\le
4K_{\mathrm{lift}}\alpha_6^{-1}\mathcal B_\bullet
\exp\!\left[
  -(1-8\delta)\frac L2\kappa\,d_{k+1}(Z)
\right].
\tag{3.1}
\]
After the designated animal-entropy spend, the rooted exponent is
\[
(1-9\delta)\frac L2\kappa.
\tag{3.2}
\]
The final source hierarchy chooses
\((1-10\delta)(L/2)\kappa=\kappa\), leaving the declared output exponent.

\begin{theorem}[Conditional fixed-regulator marked first jet]
\label{thm:marked-first-jet}
Assume the complete fixed-gas source hierarchy recorded in claims
\claim{YM-RG-020}--\claim{YM-RG-027}, including the separate ordinary
pinned ceiling, the doubled marked-amplitude condition, periodic
compatibility \(LM\mid N\), and a nonempty common shifted-branch domain.
Then one interior plaquette mark passes exactly through the finite
\(\RG\)-II algebra.  Its final connected derivative at \(t=0\) is absolutely
summable with the output polymer exponent \(\kappa\), and completed
shifted-branch numerator/denominator pairs represent the same projective
first jet without an orbit factor.
\end{theorem}

\textbf{Boundary of Theorem \ref{thm:marked-first-jet}.}
It is a first derivative at \(t=0\) in independent external coordinates.
It neither constructs source-dependent activities on a common \(t\)-disk
nor converts the derivative to the physical minimizing coarse field.

\subsection{Geometry, animals, and completed support}

For the declared closed-cube model, the literal union of a connected tuple
is the least admitted hull.  A contained-tree metric and rooted-animal count
give an explicit sufficient entropy window.  In four dimensions the
convenient threshold appearing in the ledger is
\[
\eta_{\mathrm{animal}}>64\log 8.
\tag{3.3}
\]
The final ordinary \(\RG\)-II gas is mapped into this abstract bound only
after keeping its one-sided metric comparison and species aggregation
explicit.

The shifted-branch construction is synchronized over dual numbers.
Periodic covering-space endpoint geometry removes an artificial tree-lift
hypothesis.  Because the completed coefficient is grouped by its literal
union and the Ursell and normalization factors are field-independent, the
external-\((U,J)\) dependence is confined to bonds whose geometric segments
meet the interior of that union.  On the aligned integer-wall cubulations,
the completed \(J\)-support has zero endpoint halo.

\section{Conditional physical pullbacks}

\subsection{The affine \(J\) tube}

Let \(\widehat{\mathcal C}_p^+(s,R;U,J)\) denote a completed connected
coefficient.  The compatibility hypothesis \(\HJ\) requires one raw affine
\(J\) representative whose restrictions agree across all overlapping
charts on each complete shifted branch.  Under that hypothesis, the
existing marked majorant can be rerun in a coefficientwise local
\(H^\infty\) norm:
\[
\sum_{s,R} e^{\kappa d_{k+1,s}(R)}
\sup_{\|J-\bar J\|_\infty<\Delta_J}
|\widehat{\mathcal C}_p^+(s,R;U,J)|
\le B_{\mathrm{conn}}.
\tag{4.1}
\]
Banach-line Cauchy and the exact support theorem then give
\[
\sum_{s,R} e^{\kappa d_{k+1,s}(R)}
\|D_J\widehat{\mathcal C}_p^+(s,R)\|
\le \frac{B_{\mathrm{conn}}}{\Delta_J},
\tag{4.2}
\]
with no bond-volume or branch factor.

The hypothesis \(\HJ\) is not proved by the source papers.  Equation (4.2)
is therefore a conditional completion theorem, not an unconditional
physical source derivative.

\subsection{A product \(U/J\) tube}

Raw independent-link \(U\) balls are poorly adapted to the curvature domain.
The ledger instead introduces a relative-log, background-covariant norm
controlling raw values, covariant gradients, and transported plaquette
curls.  Under the real-center hypothesis \(\Hrc\), every retained physical
center has one real/unitary representative with a uniform curvature margin.
Ordered matrix estimates then give a positive product \(U/J\) tube and full
dual derivative norms.

Neither the primary-source audit nor the finite tests prove \(\Hrc\) for
the entire completed minimizing family.  This is one of the main unresolved
uniformity assumptions.

\subsection{Support-anchored physical \(J\)}

Let \(S_{p,s,R}\) be the complete active external-\(J\) support of a
coefficient.  After composing the completed \(J\)-derivative with the
audited source kernel, the repository obtains, under its common kernel and
chart constants,
\begin{align}
&\sum_{s,R} e^{\kappa d_{k+1,s}(R)}
\sum_{y'}\mu(y')
e^{\gamma d_{\mathcal B}(S_{p,s,R},y')}
\bigl|\mathcal L^{J,\mathrm{conn}}_{p,s,R}(y')\bigr|
\notag\\
&\hspace{3em}\le
\frac{
  C_\chi c_1(\alpha_\gamma)
  \overline E_J^{\mathrm{conn}}B_{\mathrm{conn}}
}{
  \Delta_{\mathrm{tube}}
},
\qquad
0\le\gamma<\frac{\delta_0}{8}.
\tag{4.3}
\end{align}
The weight is anchored at the \emph{full active support}, not automatically
at the marked plaquette.  A stronger plaquette-rooted theorem needs an
additional endpoint or coefficient-weighted hypothesis.

\subsection{Physical \(U\), quotient synthesis, and transported curl}

For an active bond set \(I\), directions invisible to the coefficient form
a nullspace \(N_I\).  The supported derivative naturally lives in the dual
of the quotient \(X_U/N_I\), not in raw coordinate \(\ell^1\).  Exact
finite-dimensional phase duality reduces the conditional physical-\(U\)
estimate to a quotient-synthesis operator norm for the converted source
columns.

The covariant curl is not an independent mystery row.  Under the real-center
curvature margin, the ledger proves the exact four-link identity
\[
\mathcal D_{\bar U}^{\xi}a
=
D_\mu^\xi a_\nu-D_\nu^\xi a_\mu
+
\xi^{-1}
\bigl(I-\operatorname{Ad}_{d\bar U(p)}\bigr)
\bigl(a_\nu+T_\nu a_\mu\bigr).
\tag{4.4}
\]
Consequently, the raw/gradient/operator/curl norm is uniformly equivalent
to a raw-plus-covariant-gradient norm on that real-center tube.  A Lipschitz
cutoff of fixed physical width \(\rho\) yields
\[
C_{\mathrm{ext}}(\rho)
\le C_{\mathrm{eq}}
\left(1+\frac{c_{\mathrm{Ad}}^{\mathrm{RG}}}{\rho}\right).
\tag{4.5}
\]
A flat longitudinal path proves that a fixed number \(m\) of lattice layers
necessarily costs order \(1/(m\xi)\).  Thus a fixed-layer collar cannot
replace the physical-width collar uniformly as \(\xi\to0\).

\section{Why the obvious strong-norm repairs fail}

\subsection{Independent-link collars}

On the finest homogeneous patch, a complex perturbation concentrated on one
bond changes an adjacent plaquette at first order.  The small-field
curvature domain allows only \(O(\xi^2)\) deviation in the relevant
coordinates.  Therefore a full complex independent-link ball contained in
that domain has radius at most \(O(\xi^2)\).  Generic Cauchy estimates on
such a ball carry an \(O(\xi^{-2})\) loss.

This is a finite-regulator analytic-domain obstruction.  It does not say
that every structured source direction is bad.

\subsection{Diagonal scale balancing}

The source kernel bounds in the variational paper contain raw and gradient
output powers of the schematic form \(s_j^{-1}\) and \(s_j^{-2}\)
\cite{balaban1985variational}.  Weighting the chart norm by \(s_j\) and
\(s_j^2\) cancels those powers algebraically.  However, the same weights
make the full chart ball's admissible radius \(O(\xi^2)\) on the finest
patch.  Banach-line Cauchy returns the \(O(\xi^{-2})\) loss.

\begin{proposition}[Scale-renorming no-go]
\label{prop:renorm-no-go}
A diagonal raw/gradient weighting that cancels the printed
\(s_j^{-1},s_j^{-2}\) row powers does not, by itself, produce a
regulator-uniform physical-\(U\) derivative theorem through generic Cauchy
estimates.  The analytic radius collapses with the same scale power that
the row weighting removes.
\end{proposition}

A separate flat curl-free longitudinal example shows that retaining an
unweighted curl component does not evade the gradient-domain obstruction.
Proposition \ref{prop:renorm-no-go} does not rule out a range-restricted
tube, a direct weak-dual estimate, or structural cancellation.

\subsection{The joint \(U/J\) Ward quotient}

A physical coarse-source column changes both the relative \(U\) field and
the auxiliary \(J\) coordinate.  Write
\[
V_\alpha=(A_\alpha,L_{\bar U}A_\alpha).
\tag{5.1}
\]
The two summands must be paired before absolute values.  A one-dimensional
bidisc example in the ledger shows that separate absolute values can destroy
an exact Ward cancellation.

Let \(H_F\) contain coefficient-inactive directions and simultaneous
complexified gauge tangents.  Conditional on the completed coefficientwise
Ward identity, the sharp coefficient-independent first-jet constant is
\[
\bigl\|Q_{H_F}K_F^w\bigr\|_{
  \ell^\infty(\text{sources})\to E_\xi/H_F}.
\tag{5.2}
\]
Arbitrary phases in the \(\ell^\infty\) source ball forbid relying on
cancellation between distinct source columns.  Genuine cancellation must
occur within each joint \(U/J\) column or in the quotient.

The Ward premise and a regulator-uniform bound for the actual converted
synthesis are not proved.  A bounded gauge-invariant Cartan plaquette
functional gives a logical same-cell transverse countermodel under the
retained abstract hypotheses, but that countermodel is not an actual
Ba{\l}aban polymer coefficient.

\section{Exact elliptic factorization and a finite-model Stokes obstruction}

\subsection{Measure-normalized columns}

Let \(c\) be a physical coarse-source coordinate with source measure
\(\mu_c\).  Use the density convention
\[
D_B\mathcal H[h]=\sum_c\mu_cA_ch_c.
\tag{6.1}
\]
The coordinate column constrained by differentiation is
\[
\kappa_c=\mu_cA_c.
\tag{6.2}
\]
At a fixed source center \(\bar B\), define the differentiated averaging
and projected-Landau operators \(C_{\bar B}\) and \(S_{\bar B}\).  Under
the repository density convention (6.1), differentiating variational-paper
Eqs.~(20)--(21) gives \cite{balaban1985variational}
\[
C_{\bar B}\kappa_c=e_c,
\qquad
S_{\bar B}\kappa_c=0.
\tag{6.3}
\]
Applying (6.3) to the density \(A_c\) instead of the coordinate column
\(\kappa_c\) loses a source-measure factor.

\begin{proposition}[Affine constrained range]
\label{prop:affine-range}
Let \(C:X\to Y\) and \(S:X\to Z\) be maps between finite-dimensional Hilbert
spaces, and assume \(C|_{\ker S}\) is onto.  If
\(\mathcal Z=\ker C\cap\ker S\) and \(R_C:Y\to\ker S\) is any right inverse,
then all solutions of \(C\kappa=e\), \(S\kappa=0\) are
\[
\kappa=R_Ce+\mathcal Z.
\tag{6.4}
\]
A detector \(\ell\in X^*\) is fixed by these constraints exactly when
\[
\ell|_{\mathcal Z}=0
\quad\Longleftrightarrow\quad
\ell\in\operatorname{ran}C^*+\operatorname{ran}S^*.
\tag{6.5}
\]
\end{proposition}

Thus averaging and projected Landau select an affine slice, not a unique
column.  The variational Hessian supplies the additional selection.

\subsection{The exact differentiated factorization}

The variational source uses three different constrained lifts
\cite{balaban1985variational}: \(H\) is the original lift used in the
nonlinear restoration \(HD(A')\), \(H_1\) is the improved lift in
Eqs.~(174)--(177), and \(H_0\) is the alternative lift in
Eqs.~(179)--(190).  They must not be identified with one another or with
the later Landau-to-axial gauge map.  Put
\[
A_B=\mathcal A_0(B)+H_0B,
\qquad
T_B=D_BA_B,
\qquad
W_B=V''(A_B).
\tag{6.6}
\]
With
\(\widetilde G=(\Delta_a-\Delta^{(2)})^{-1}\), the differentiated equation
in the variational paper is
\[
\frac{\delta\mathcal A_0}{\delta B}
+\widetilde G\,W_BT_B
=\widetilde G\Delta^{(2)}H_0.
\tag{6.7}
\]
Since \(\delta_B\mathcal A_0=T_B-H_0\), multiplication by
\(\widetilde G^{-1}\) gives
\[
\bigl(\Delta_a-\Delta^{(2)}+W_B\bigr)T_B
=\Delta_aH_0.
\tag{6.8}
\]
The background derivative is therefore
\[
\boxed{
K_B=D_B\mathcal H(B)
=
\bigl(I-HD'(A_B)\bigr)
\bigl(\Delta_a-\Delta^{(2)}+V''(A_B)\bigr)^{-1}
\Delta_aH_0.
}
\tag{6.9}
\]
Equation (129) of the source, together with
\(G=\Delta_a^{-1}\), implies
\[
\Delta_aH_0
=
Q^*(QGQ^*)^{-1}
\bigl(L^{j(\cdot)}\eta\bigr)^{-1}.
\tag{6.10}
\]
The order-\(-2\) inverse in (6.9) is fed by an order-\(+2\) source.  It
cannot be read in isolation as a two-power smoothing theorem.  The printed
expansion also yields the exact consequence
\[
K_0=D_B\mathcal H(0)=H_1.
\tag{6.11}
\]
Equations (6.8) and (6.9) are derived from the printed differentiated
Eqs.~(182)--(183); (6.10) is derived from printed Eq.~(129) and the
preceding definition \(G=\Delta_a^{-1}\); and (6.11) is the derivative
consequence of the printed expansion (177)
\cite{balaban1985variational}.  The Stokes calculation below has a
deliberately weaker status.

\subsection{A source-normalized finite periodic model}

Take
\[
\Lambda_f=(\mathbb Z/4\mathbb Z)^2
\times(\mathbb Z/2\mathbb Z)^2,
\qquad
\Lambda_c=(\mathbb Z/2\mathbb Z)^2
\times(\mathbb Z/1\mathbb Z)^2,
\tag{6.12}
\]
and set \(L=2\).  Restrict to a fixed Cartan generator \(T\in\mathfrak{su}(2)\)
and fields constant in the last two coordinates.  Write
\(a=\eta\kappa\).  The finite model retains the flat-Cartan \(Q_0,Q'_0\)
and projected-Landau normalizations:
\begin{align}
(Ca)_1(X,Y)
&=\frac14\sum_{r\in\{0,1\}^2}\sum_{t=0}^1
a_1\bigl(2(X,Y)+r+te_1\bigr),
\notag\\
(Ca)_2(X,Y)
&=\frac14\sum_{r\in\{0,1\}^2}\sum_{t=0}^1
a_2\bigl(2(X,Y)+r+te_2\bigr).
\tag{6.13}
\end{align}
The scalar block map is the full four-dimensional mean; on invariant fields
it restricts to the \(2\times2\) mean.  The model Landau projector is the
orthogonal projection onto \(\Delta\ker P_4\).

\begin{remark}[Source boundary]
The averaging and gauge-fixing papers support the normalizations used in
(6.13) \cite{balaban1985averaging,balaban1985gaugefixing}.  They do not,
in the material audited here, prove that the anisotropic finite periodic
model (6.12) is an embedded patch of the complete multiscale \(H_1\)
construction, its full \(Q'\) nullspace, or its boundary conditions.
\end{remark}

\subsection{Exact representative and discrete Stokes}

Let \(e_c\) be one positive coarse \(1\)-bond at \((X,Y)=(0,0)\).
The following rational field is extended constantly in the spectator
coordinates, with \(a_2=a_3=a_4=0\):
\[
a_1=
\begin{pmatrix}
\frac14&\frac34&\frac14&-\frac14\\
\frac14&\frac34&\frac14&-\frac14\\
0&0&0&0\\
0&0&0&0
\end{pmatrix}.
\tag{6.14}
\]
Direct substitution gives \(Ca=e_c\).  If \(g=d^*a\), then
\[
g=
\begin{pmatrix}
-\frac12&-\frac12&\frac12&\frac12\\
-\frac12&-\frac12&\frac12&\frac12\\
0&0&0&0\\
0&0&0&0
\end{pmatrix}.
\tag{6.15}
\]
The field \(\Delta g\) is constant on each full \(2^4\) block.  Hence
\(\Delta g\in\operatorname{ran}P_4^*=(\ker P_4)^\perp\), so
\[
g\perp\Delta\ker P_4,
\qquad
Rd^*a=0.
\tag{6.16}
\]
Thus the model's localized right-inverse/Landau slice is nonempty.

The bond map is a cochain map:
\[
d_cC=C_2d,
\tag{6.17}
\]
where a selected row of \(C\) has \(\ell^1\)-mass \(2\) and a selected row
of \(C_2\) has \(\ell^1\)-mass \(4\).  Since
\((d_ce_c)(0,0)=1\), every model field with \(Ca=e_c\) obeys
\[
\|a\|_\infty\ge\frac12,
\qquad
\|da\|_\infty\ge\frac14,
\qquad
\|\delta a\|_\infty\ge\frac18.
\tag{6.18}
\]
With \(s=2\eta\) and \(\kappa=\eta^{-1}a\), this becomes:

\begin{theorem}[Finite-model coarse-curl obstruction]
\label{thm:stokes}
Every right inverse of the localized source in the finite model
(6.12)--(6.13), with or without the projected-Landau condition, satisfies
\[
\boxed{
\|\kappa\|_\infty\ge s^{-1},
\qquad
\|\nabla^\eta\kappa\|_\infty\ge\frac12s^{-2},
\qquad
\|d^\eta\kappa\|_\infty\ge s^{-2}.
}
\tag{6.19}
\]
The representative (6.14) has matching powers:
\[
\|\kappa\|_\infty\le\frac32s^{-1},
\qquad
\|\nabla^\eta\kappa\|_\infty,\,
\|d^\eta\kappa\|_\infty\le3s^{-2}.
\tag{6.20}
\]
\end{theorem}

Theorem \ref{thm:stokes} is an exact rational finite-dimensional theorem.
It proves that right-inverse normalization and projected-Landau
orthogonality alone do not force a two-power strong-norm gain in this model.
It does \emph{not} prove that the actual \(H_1\) column contains this model,
that the full Hessian selector retains the curl mode, or that a completed
physical coefficient detects it.

\section{Research handoff}

\subsection{The next decision, in order}

The most informative next work is not another generic strong-norm estimate.
It is the following decision tree.

\begin{enumerate}
  \item \textbf{Embedding test.}
  Prove or disprove that the homogeneous Cartan finite model can be embedded
  into the full multiscale \(H_1\) domain with its complete \(Q'\) nullspace,
  domain decomposition, boundary conditions, and Hessian selection.

  \item \textbf{If the curl survives, test the completed weak dual.}
  For each completed coefficient, evaluate the actual within-column
  functional
  \[
  D_U\widehat{\mathcal C}[A_\alpha]
  +
  D_J\widehat{\mathcal C}[L_{\bar U}A_\alpha]
  \tag{7.1}
  \]
  on the surviving coarse-curl range after the simultaneous Ward quotient.
  Either prove annihilation/cancellation or construct a source-faithful
  completed coefficient that detects the mode.

  \item \textbf{If strong columns remain incompatible, change the analytic
  object honestly.}
  Formulate a range-restricted tube, a direct weak-quotient-dual estimate,
  or a physically justified curl-free, smoothed, weighted, or separated
  source class.  Do not silently change the intended plaquette observable.

  \item \textbf{Complete the existing first-jet prerequisites.}
  Prove \(\Hrc\), the common converted raw/gradient rows, the completed Ward
  identity, the active-label metric halo when required, and the compatible
  \(J\)-representative atlas.

  \item \textbf{Only then build the nonzero-source and iterative theory.}
  Construct source-dependent polymer activities on a common disk, add large
  fields, prove unit translations and projection, iterate the \(\RG\), and
  begin regulator removal.
\end{enumerate}

\subsection{Assumption ledger}

\begin{longtable}{>{\raggedright\arraybackslash}p{0.23\textwidth}
                  >{\raggedright\arraybackslash}p{0.30\textwidth}
                  >{\raggedright\arraybackslash}p{0.37\textwidth}}
\toprule
\textbf{Item} & \textbf{What exists} & \textbf{What remains} \\
\midrule
\endfirsthead
\toprule
\textbf{Item} & \textbf{What exists} & \textbf{What remains} \\
\midrule
\endhead
Selected branch &
Finite normalized law, local holomorphy, gauge-covariant push-forward &
Equality with the unrestricted raw transform; global chart compatibility \\

Marked first jet &
Exact dual-number routing and conditional fixed-gas convergence at \(t=0\) &
Common nonzero-source polymer disk and second source jet \\

Completed support &
Literal-union external support and zero \(J\)-endpoint halo on aligned cells &
Unit-translation theorem and any stronger marked-root metric halo \\

Affine \(J\) tube &
Conditional coefficientwise \(H^\infty\) rerun and dual bound &
\(\HJ\) compatible representative atlas \\

Product \(U/J\) tube &
Explicit norm and radius under a real-center curvature margin &
\(\Hrc\) for the full physical family \\

Source kernel &
Audited row bounds, exact affine range, exact elliptic factorization &
Complete chart/restoration conversion and uniform physical synthesis \\

Stokes mode &
Exact source-normalized finite-model lower bound &
Embedding into full multiscale \(H_1\) and Hessian-selected survival \\

Ward quotient &
Exact finite-dimensional quotient/phase theorem &
Completed coefficientwise Ward identity and actual uniform quotient norm \\

Large fields and iteration &
Source-free route-level literature context &
Source-dependent large-field operation, iteration, projection, running flow \\

Continuum and infrared &
Only a roadmap and conditional reconstruction lemmas &
Every construction, nontriviality, OS, infinite-volume, decay, and gap gate \\
\bottomrule
\end{longtable}

\subsection{Clay obligation matrix}

\begin{center}
\begin{tabularx}{\textwidth}{>{\raggedright\arraybackslash}p{0.28\textwidth}
                             >{\raggedright\arraybackslash}X}
\toprule
\textbf{Obligation} & \textbf{State at this handoff} \\
\midrule
Every compact simple group & Open; the calculations are primarily \(SU(2)\). \\
Continuum theory on \(\mathbb R^4\) & Open; only finite regulators appear. \\
Infinite volume & Open. \\
Nontriviality & Open. \\
Local continuum observables and UV behavior & Open. \\
OS/Wightman-strength axioms and reconstruction & Open. \\
Finite positive physical mass gap & Open; no infrared bridge exists. \\
Independent expert review & None recorded. \\
\bottomrule
\end{tabularx}
\end{center}

\section{Reproducibility and source boundary}

\subsection{Repository map}

The authoritative machine-readable state is in
\path{STATUS.json}, \path{CLAIMS.json}, and
\path{audit/objections.json}.  The principal research narrative is in
\path{research/LOG.md}.  The most relevant notes are grouped below.

\begin{center}
\begin{tabularx}{\textwidth}{>{\raggedright\arraybackslash}p{0.22\textwidth}
                             >{\raggedright\arraybackslash}X}
\toprule
\textbf{Notes/claims} & \textbf{Role} \\
\midrule
0020--0022 &
Distinguished-vertex Ursell identity, exact marked Section-2
factorization, and whole-integrand Cauchy. \\
0023--0024 &
Literal-union geometry, rooted animals, and explicit ordinary pinned
Koteck\'y--Preiss interface. \\
0025--0027 &
Source-faithful conditioning, positive marked resummation, and shifted
first-jet synchronization. \\
0028--0029 &
Periodic endpoint bridge and completed external support. \\
0030--0032 &
Conditional affine/product tubes and support-anchored physical \(J\). \\
0033--0034 &
Physical-\(U\) quotient synthesis, transported-curl reduction, and
physical-width extension. \\
0035 &
Diagonal scale-renorming no-go. \\
0036 &
Joint Ward quotient, vertical scale trade, and an abstract transverse
countermodel. \\
0037 &
Exact elliptic factorization plus the source-normalized finite-model
coarse-curl Stokes obstruction. \\
\bottomrule
\end{tabularx}
\end{center}

The focused finite-model regression is
\path{tests/test_elliptic_range_factorization_stokes_obstruction.py}.
Repository-wide checks are run with
\begin{verbatim}
python scripts/verify_repository.py
python -m unittest discover -s tests -v
\end{verbatim}
The frozen handoff snapshot is checkpoint \claim{YM-RG-037} on branch
\texttt{codex/first-source-audit}, with 141 passing tests and public draft
\href{https://github.com/YesterdaysLemon/yang-mills-research/pull/5}
{pull request 5}.  The final immutable commit identifier is recorded in the
repository history and pull-request head rather than embedded circularly in
this artifact.
These commands certify neither the primary literature nor the deep
mathematics.  They are guards against algebraic drift and status inflation.

\subsection{Primary-source audit boundary}

The retained local audits visually checked the relevant pages of:
\begin{itemize}
  \item Ba{\l}aban's averaging paper for the linear and nonlinear block-map
  normalizations \cite{balaban1985averaging};
  \item the gauge-fixing paper for the projected Landau operator
  \cite{balaban1985gaugefixing};
  \item the variational paper for the three distinct lifts, the background
  expansion, and the differentiated equations
  \cite{balaban1985variational};
  \item \emph{Propagators II} for the multiscale distance/decay transfer
  \cite{balaban1984propagatorsII};
  \item \(\RG\) I and II for the selected small-field and cluster scaffolds
  \cite{balaban1987rgi,balaban1988rgii}.
\end{itemize}
Every imported-versus-derived distinction should be preserved.  In
particular:
\begin{itemize}
  \item \(H,H_0,H_1\) are distinct operators;
  \item (6.9) and (6.10) are algebraic consequences, not separately printed
  displays;
  \item (6.11) is the derivative consequence of a printed expansion;
  \item Theorem \ref{thm:stokes} is repository finite-model algebra, not a
  theorem printed by Ba{\l}aban and not a proved \(H_1\) embedding.
\end{itemize}

\section{Conclusion}

The project did not solve the \YM{} existence and mass gap problem.  It did
produce a useful finite-regulator map of one local-observable bottleneck.
The one-mark cluster algebra is substantially organized; exact support and
quotient structures are known; and several seductive norm-based repairs
have concrete countermodels.  The differentiated variational equations have
an exact elliptic factorization, while the final Stokes model shows why the
normalized range constraints alone should not be mistaken for strong
smoothing.

The most important negative discovery is methodological: a Green operator,
a weighted norm, a Ward slogan, or a projected Landau condition cannot be
read separately from its source normalization, analytic radius, coefficient
dual, and domain.  The most important positive handoff is a sharp next
experiment: determine whether the full multiscale Hessian-selected
\(H_1\) range contains the finite-model curl mode, and if it does, test the
completed simultaneous \(U/J\) weak dual on that range before investing in
further strong-column estimates.

\appendix

\section{Exact nonclaims}

For avoidance of doubt, this paper does not establish any of the following:
\begin{enumerate}
  \item equality between the selected small-field branch and the
  unrestricted raw block transform;
  \item source-dependent polymer activities on a common nonzero source disk;
  \item a common real-center theorem \(\Hrc\) for the completed physical
  family;
  \item converted raw/gradient source rows with regulator-uniform constants;
  \item the completed coefficientwise simultaneous \(U/J\) Ward identity;
  \item embedding of the finite Stokes model into the full \(H_1\) range;
  \item a source-dependent large-field theorem;
  \item \(\RG\) iteration, marginal projection, or control of the running
  coupling through a dynamical mass scale;
  \item continuum or infinite-volume existence, nontriviality, Euclidean
  covariance, OS reconstruction, physical-unit exponential decay, or a
  mass gap;
  \item extension from the initial \(SU(2)\) calculations to every compact
  simple group.
\end{enumerate}

\section{Future-agent falsification checklist}

A future agent should reject any proposed continuation that:
\begin{itemize}
  \item applies the differentiated right-inverse identity to \(A_c\) rather
  than \(\kappa_c=\mu_cA_c\);
  \item identifies \(H,H_0,H_1\), or calls the nonlinear restoration a gauge
  transformation;
  \item reads the inverse in (6.9) as smoothing while deleting its
  \(\Delta_aH_0\) source;
  \item replaces projected Landau \(RD^*K=0\) by full Landau \(D^*K=0\);
  \item calls the anisotropic finite periodic model an actual multiscale
  source patch without proving the embedding;
  \item takes separate absolute values of the \(U\) and induced-\(J\)
  summands before using the Ward quotient;
  \item relies on cancellation between source columns despite arbitrary
  source phases;
  \item spends output decay without checking the same-layer same-cell
  diagonal;
  \item hides an \(O(\xi^2)\) analytic radius inside a weighted norm constant;
  \item promotes a finite-regulator, conditional, or model statement to a
  continuum construction or mass gap.
\end{itemize}

\begin{thebibliography}{99}
\raggedright
\small

\bibitem{jaffewitten}
A.~Jaffe and E.~Witten,
\emph{Quantum Yang--Mills Theory},
official Clay Mathematics Institute problem description,
\url{https://www.claymath.org/wp-content/uploads/2022/06/yangmills.pdf}.

\bibitem{wilson1974}
K.~G.~Wilson,
\emph{Confinement of quarks},
Physical Review D \textbf{10} (1974), 2445,
\href{https://doi.org/10.1103/PhysRevD.10.2445}
{doi:10.1103/PhysRevD.10.2445}.

\bibitem{os1973}
K.~Osterwalder and R.~Schrader,
\emph{Axioms for Euclidean Green's functions},
Communications in Mathematical Physics (1973),
\href{https://doi.org/10.1007/BF01645738}{doi:10.1007/BF01645738}.

\bibitem{os1975}
K.~Osterwalder and R.~Schrader,
\emph{Axioms for Euclidean Green's functions II},
Communications in Mathematical Physics (1975),
\href{https://doi.org/10.1007/BF01608978}{doi:10.1007/BF01608978}.

\bibitem{osterwalderseiler}
K.~Osterwalder and E.~Seiler,
\emph{Gauge field theories on a lattice},
Annals of Physics (1978),
\href{https://doi.org/10.1016/0003-4916(78)90039-8}
{doi:10.1016/0003-4916(78)90039-8}.

\bibitem{luscher}
M.~L\"uscher,
\emph{Construction of a self-adjoint, strictly positive transfer matrix
for Euclidean lattice gauge theories},
Communications in Mathematical Physics (1977),
\href{https://doi.org/10.1007/BF01614090}{doi:10.1007/BF01614090}.

\bibitem{balaban1984propagatorsII}
T.~Ba{\l}aban,
\emph{Propagators and renormalization transformations for lattice gauge
theories II},
Communications in Mathematical Physics \textbf{96} (1984), 223--250,
\href{https://doi.org/10.1007/BF01240221}{doi:10.1007/BF01240221}.

\bibitem{balaban1985averaging}
T.~Ba{\l}aban,
\emph{Averaging operations for lattice gauge theories},
Communications in Mathematical Physics \textbf{98} (1985), 17--51,
\href{https://doi.org/10.1007/BF01211042}{doi:10.1007/BF01211042}.

\bibitem{balaban1985gaugefixing}
T.~Ba{\l}aban,
\emph{Spaces of regular gauge field configurations on a lattice and gauge
fixing conditions},
Communications in Mathematical Physics \textbf{99} (1985), 75--102,
\href{https://doi.org/10.1007/BF01466594}{doi:10.1007/BF01466594}.

\bibitem{balaban1985backgroundpropagators}
T.~Ba{\l}aban,
\emph{Propagators for lattice gauge theories in a background field},
Communications in Mathematical Physics \textbf{99} (1985), 389--434,
\href{https://doi.org/10.1007/BF01240355}{doi:10.1007/BF01240355}.

\bibitem{balaban1985variational}
T.~Ba{\l}aban,
\emph{The variational problem and background fields in renormalization
group method for lattice gauge theories},
Communications in Mathematical Physics \textbf{102} (1985), 277--309,
\href{https://doi.org/10.1007/BF01229381}{doi:10.1007/BF01229381}.

\bibitem{balaban1987rgi}
T.~Ba{\l}aban,
\emph{Renormalization group approach to lattice gauge field theories I:
Generation of effective actions in a small field approximation and a
coupling constant renormalization in four dimensions},
Communications in Mathematical Physics \textbf{109} (1987), 249--301,
\href{https://doi.org/10.1007/BF01215223}{doi:10.1007/BF01215223}.

\bibitem{balaban1988rgii}
T.~Ba{\l}aban,
\emph{Renormalization group approach to lattice gauge field theories II:
Cluster expansions},
Communications in Mathematical Physics \textbf{116} (1988), 1--22,
\href{https://doi.org/10.1007/BF01239022}{doi:10.1007/BF01239022}.

\bibitem{koteckypreiss}
R.~Koteck\'y and D.~Preiss,
\emph{Cluster expansion for abstract polymer models},
Communications in Mathematical Physics \textbf{103} (1986),
\href{https://doi.org/10.1007/BF01211762}{doi:10.1007/BF01211762}.

\end{thebibliography}

\end{document}
